% ============================================================
%  R. Nielsen, "Om theoretisk og praktisk Erkjendelse, nogle Bemærkninger".
%  Nordisk Universitets-Tidsskrift, Tillægshefte til tiende Aargang,
%  udg. af A. Ingerslev. Kjöbenhavn 1866 (trykt hos S. Trier), pp. 63–77.
%
%  TRANSCRIPTION (Danish) — from the Google Books scan.
%  Scan: ~/bibliotek/Nielsen, Rasmus/1864-1866-nordisk-universitets-tidskrift-10.pdf
%        (Google Books VEdQAQAAMAAJ, 708 PDF pp., the whole of vol. 10 with the
%        1866 Tillægshefte bound in at the end; sha256 in SCANS.tsv.)
%
%  PAGE MAP (verified 2026-09-21: folios read off the text layer of every page
%  PDF 600–656, endpoints eye-confirmed on the images):
%        printed 63–77  =  PDF + 574   (PDF 637–651), uniform, no gaps, no dupes
%        The Tillægshefte has its own pagination. Folio and running head are
%        suppressed on p. 63 (the article's head); PDF 652 is printed 78, the
%        head of the Nekrologer (J. G. Forchhammer), folio likewise suppressed.
%
%  TYPEFACE: ANTIQUA throughout, with NO ITALIC anywhere — Latin and German
%        phrases are set in the same roman. The head: title in large roman
%        capitals over two lines, subtitle in smaller bold capitals, byline
%        bold, a short centred rule. The Greek of p. 74 is in a sloping Greek
%        fount.
%
%  CONVENTIONS: diplomatic transcription; 1866 orthography as printed;
%  printer's errors transcribed as printed and logged in % comments at the
%  site, never silently corrected. Printed line breaks are kept; words broken
%  at a line end are closed up with \-% .
%   ö     — every o-diacritic in the article is ö (including Överst); there is
%        no ø. aa throughout, never å.
%   EMPHASIS is LETTERSPACING (Sperrsatz) ONLY, rendered \emph{}: six places,
%        on pp. 64, 66, 67 and 76.
%   QUOTES „…“ as printed; the count balances (16/16).
%   SPACED PUNCTUATION before , ; : is justification, not copy, and is closed
%        up.
%   GREEK (Rom. 1:19–20, p. 74) typed in Unicode with the accents and
%        breathings the page shows; the fount's variant sigma, theta and kappa
%        normalised to σ θ κ. Two printer's readings kept: ἣ (rough breathing,
%        grave) and αυτοῦ (no breathing).
%   NOT TRANSCRIBED: running heads, folios, sheet signatures, the Google
%        watermark.
%
%  FOOTNOTE: one, p. 67, marked *) (Heiberg, Prosaiske Skrifter I, S. 465).
%
%  STRUCTURE: none — no sections or numbering. The article ends on p. 77
%        („… et intellectuelt Amphibium.“) with a short centred rule; no
%        signature, no date.
%
%  DEFECTS AND DOUBTFUL READINGS (all logged at the site): p. 68 an ellipsis
%        of three close-set points; p. 75 „Rom. 11“ with an ink-blotted first
%        1; p. 76 „Nogen“ with a stray grave-like mark on the o; p. 77
%        „Villen“ (twice) as printed.
%
%  Word-accuracy: every batch diffed against the embedded Google text layer
%  before splicing (ocrdiff.py). 1 batch, 1 candidate, 0 transcription errors
%  corrected, 0 unresolved. Cross-check 2026-09-21: a whole-article word
%  alignment against the same layer matched 4026 of 4071 words (98.9%); every
%  residual difference inspected was an OCR split or misreading. All 15 pages
%  were also rendered and read at the image by the transcribing agent; one
%  emphasis call (p. 76) spot-checked at the image by the caller.
%  See TRANSCRIPTION-PLAYBOOK.md §3 step 4.
%
%  STATUS: COMPLETE. pp. 63–77 transcribed 2026-09-21 in one batch.
% ============================================================
\documentclass[12pt,a4paper]{article}
\usepackage[utf8]{inputenc}
\usepackage[T1]{fontenc}
\usepackage[danish]{babel}
\usepackage{textalpha}   % Greek typed directly (Rom. 1:20 on p. 74)
\usepackage[protrusion=true,expansion=true]{microtype}
\usepackage[margin=1.3in]{geometry}
\usepackage{setspace}
\usepackage{libertinus}
\usepackage{libertinust1math}
\usepackage{fancyhdr}
\usepackage[colorlinks=true,linkcolor=black,urlcolor=black,
  pdftitle={Om theoretisk og praktisk Erkjendelse, nogle Bemærkninger},
  pdfauthor={R. Nielsen}]{hyperref}
\setlength{\headheight}{15pt}
\pagestyle{fancy}
\fancyhf{}
\fancyhead[L]{\textit{Om theoretisk og praktisk Erkjendelse}}
\fancyhead[R]{\thepage}
\setlength{\parindent}{1.5em}
\setstretch{1.3}
\renewcommand{\thefootnote}{*)}
\usepackage{marginnote}
\newcommand{\opage}[1]{\marginnote{\footnotesize\textit{[#1]}}}
\newcommand{\apage}[1]{\marginnote{\footnotesize\textit{[#1]}}}
\begin{document}

% --- p. 63 ---
\opage{63} % Head of the article (PDF 637). Folio and running head suppressed.
% Title in large roman capitals over two lines; subtitle in smaller bold
% capitals; byline in a bold (Clarendon-like) face; a short centred rule below.
% The first letter of the text is a slightly enlarged R, not a drop cap.
\begin{center}
  {\large OM THEORETISK OG PRAKTISK\\
  ERKJENDELSE,}

  \smallskip
  {\small\textbf{NOGLE BEMÆRKNINGER}}

  \smallskip
  \textbf{af R.\ Nielsen.}
\end{center}

\begin{center}\rule{0.2\linewidth}{0.4pt}\end{center}
\medskip

Religion og Videnskab ere store Magter i Aandens Ver\-%
den. At disse ideale Magter ofte og hæftig have bekriget
hinanden, kan ikke undre os, naar vi betænke, at de i
Oprindelse og Væsen ere saa ueensartede, og at enhver af
dem for sig, naar det gjælder Ledningen af Menneske\-%
slægtens Opdragelse og Dannelse, gjör Fordring paa
Hegemoniet. Religionen forudsætter en Aabenbaring,
Aabenbaringen et Under; men Videnskaben negter Un\-%
deret, forkaster Aabenbaringen og kommer derved til at
krænke Religionen. Hvor fordærvelig Striden er, og
hvor heldbringende det vilde være, ikke alene for Menig\-%
hedslivet, men for Folkeliv og Culturudvikling i det Hele,
dersom Misforstaaelsen, den gamle Misforstaaelse, som
ved at stille de to Aandsmagter fjendtlig imod hinanden
gjennem Aarhundreder har medvirket til at splitte og
svække de ædleste Kræfter og forvirre Opfattelsen af det,
som er og maa være alle Bestræbelsers Endemaal, om\-%
sider lod sig hæve, vil Enhver, der alvorlig har tænkt
over Sagen, sikkerlig indrömme.

Skulde det nu virkelig være muligt, at en saa dybt
indgroet Misforstaaelse som den imellem Tro og Viden,
% --- p. 64 ---
\opage{64}Religion og Videnskab, nogensinde fra Grunden af kunde
lade sig hæve?

Det synes jo næsten umuligt, naar man ret betænker,
hvormeget her kræves til god Forstaaelse. Til god
Forstaaelse imellem Tro og Viden er det nemlig ikke
nok, at Parterne omsider bekvemme sig til at tolerere
hinanden, saa at Nogle i Samfundet faae Lov at dyrke
Videnskaben, uden at bekymre sig om Troeslivet, me\-%
dens det store Fleertal, uden Indsigt i Naturens og Viden\-%
skabens Love, troskyldig antager, hvad Theologer og
Præster efter Tid og Leilighed finde tjenligt at forklare
de Troende om Troens Mirakler.

En saadan Overeenskomst medförer kun en charak\-%
teerlös Mellemtilstand. Men en charakteerlös Mellem\-%
tilstand har noget for Aandslivet sörgelig Afslövende; det
Velgjörende er, at der under en væbnet Neutralitet kan
ulme en Gnist, der, inden man veed deraf, slaaer ud i Lue.

Til god Forstaaelse mellem Religion og Videnskab
hörer, at de to Magter væsentlig maae kunne anerkjende
hinanden og under gjensidig Anerkjendelse enes om at
virke sammen til eet og samme höieste Maal. Men for
at dette kan skee, maae de to Principer, Troen og Viden,
lade sig \emph{forene i een Bevidsthed}. Den Troende
maa da, uden at tage Skade paa sin Tro, kunne blive
en Vidende, den Vidende, uden at fornegte eller halvere
sin Viden, være en Troende; det maa for den alsidig
Dannede, med andre Ord, være muligt at give Troen
hvad Troens er, og Viden hvad Videns er.

Er Fordringen ligefrem, frit og aabent udtalt, da
ville Indsigelserne trænge sig frem med aandelig Natur\-%
nödvendighed. „Skulde Troens og Videns absolut ueens\-%
artede Principer virkelig lade sig forene i een Bevidsthed,
da maa det samme Menneske jo som troende antage,
hvad han som vidende fornegter, Aabenbaringens Under;
han maa altsaa paa een Gang baade anerkjende og dog
ikke anerkjende det Selvsamme; men det er en utaalelig
% --- p. 65 ---
\opage{65}Modsigelse. Den, der trods Modsigelsen vil forsöge paa at
löse en saadan Opgave, maa da begynde med at halvere sin
Bevidsthed. Ligesom hiin Konge, der smilede med det
halve Ansigt og truede med det halve, i en vis Forstand
viste sig som en Mand med to Ansigter, saaledes maatte den
Bevidsthed, der skulde forene de to ueensartede Principer,
være vidende med den ene Halvdeel og troende med den
anden, have en Sandhed i den ene, en anden i den anden,
og saaledes vise sig med to Sandheder (duæ veritates).
Og hvad kunde det saa hjælpe? Saalænge de to Sand\-%
heder kun forenes i een Bevidsthed saaledes, at de som
ved en Slags intellectuel Mellemvæg holdes afsondrede
hver i sit Rum, er den Modsætning, som tidligere vakte
Strid og Splid i det Udvortes, jo blot flyttet ind i det
Indvortes. Brister Mellemvæggen, og den maa vel briste,
naar Aanden engang afryster det Mechaniske, saa er der
jo Fare for, at de ueensartede Sandheder, idet de forenes,
brænde sammen med en Explosion, der sprænger For\-%
stand og Hjerne.“

Ved saaledes at blotte den formeentlige Modsigelse,
have vi opnaaet at see Problemet blive til. Men at et
Problem maa blive til, er jo den förste Betingelse for,
at det skal kunne löses. Kan det da löses? Det gjælder
et Forsög.

Den „store, gruelige, forfærdelige Modsigelse“, man
har villet finde i den Omstændighed, at Bevidstheden ved
at anerkjende to absolut ueensartede Principer skulde
spalte sig paa en saadan Maade, at den tabte sin Een\-%
hed og dermed sin Energie, har, naar vi see nöiere til,
ikke meget at betyde. Enhver, der har indövet saamegen
Dialektik, at han har lært at indsee, hvad det er, hvor\-%
paa det i Forholdet mellem et Absolut og et Relativt
væsentlig kommer an, vil uden Vanskelighed overbevise
sig om, at Principer, der ere absolut ueensartede, meget
vel kunne være i indbyrdes Relation og i denne Relation
vise sig relative. Enhver, der har tilegnet sig saamegen
% --- p. 66 ---
\opage{66}Psychologie, at han kjender de förste Grundtræk til Læren
om Bevidsthedens Eenhed under dens gjennemgaaende
Selvfordoblinger, vil ved nogen Eftertanke kunne sige sig
selv, at der i Aandslivet gives Modsætninger, dybe, stærke,
til Dualisme grændsende Modsætninger, der dog ere saa
langt fra at svække, at de tvertimod forhöie og bekræfte
Selvbevidsthedens Eenhed.

Til de almindeligste og derfor meest bekjendte For\-%
mer for Fordoblinger af Aandslivet hörer uden Tvivl
Modsætningen mellem Theorie og Praxis, mellem theore\-%
tisk og praktisk Erkjendelse.

At Theorie og Praxis maae kunne staae i Relation,
ja i væsentlig Relation til hinanden, skal vist Ingen
negte; men desuagtet have de dog hver sit Princip og
ere forsaavidt ueensartede. Tör en theoretisk Begrun\-%
delse ikke indblande praktiske Hensyn, saa kan den
virkelige Praxis heller ikke ligefrem afledes af Theorien.
Vilde Nogen spörge: hvad er i Bevidsthedslivet det Op\-%
rindelige, den theoretiske eller den praktiske Erkjen\-%
delse? hvo skulde da betænke sig paa et Svar.

Hvor gjennemgaaende Modsætningen mellem theore\-%
tisk og praktisk Erkjendelse er, kan da allerbedst sees
af Bevidsthedens tilsvarende Forvandling. Den theore\-%
tiske Bevidsthed er upersonlig, den praktiske er person\-%
lig; hiin er contemplativ, uinteresseret, uden Selvhensyn
fortabt i det Almene, bevæget af Nödvendigheden, glem\-%
mende sig selv i Begrundelsen af det Objective; denne
derimod er interesseret, i Conflict med Existensen, be\-%
væget af Frygt og Haab, Attraa og Forventning, endog
i sin Selvfornegtelse fuld af Selvhensyn. Den theoretiske
Bevidsthed er væsentlig \emph{tænkende}, den praktiske væ\-%
sentlig \emph{villende}.

At der hörer Villie til for at udvikle Tanken, og
Tanke til for at anvende Villien, er begribeligt; Tanken
og Villien, Bevidsthedens theoretiske og praktiske Fac\-%
torer, befinde sig ganske vist i indbyrdes Relation og
% --- p. 67 ---
\opage{67}ere i deres Forskjellighed forsaavidt relative; men Tan\-%
ken er derfor ikke Villie, Villien ikke Tanke; Theorien
er ikke Praxis, Praxis ikke Theorie: den theoretiske
Bevidsthed, der giver sig hen i Tanker og Tænkning,
lever i en anden Region end den praktiske, der con\-%
centrerer sig i Forsæt og Villie.

Det er en Misforstaaelse, en for Religion, Sædelig\-%
hed og al höiere Charakteerudvikling dödbringende Mis\-%
forstaaelse, at gjöre Villien til en Modification af Tænk\-%
ningen, en Misforstaaelse, som, hvad den nyere Philosophie
angaaer, kan siges at begynde med Spinoza (voluntas
et intellectus unum et idem sunt), for at fuldendes ved
Hegel. I den hegelske Philosophie er det das Denken
und das Denken und das allgemeine Denken, hvorpaa
det först og sidst kommer an. Fra dette Synspunkt har
da ogsaa I.\ L.\ Heiberg i sit „Indledningsforedrag til det
logiske Cursus“\footnote{% printed mark *)
Prosaiske Skrifter. Förste Bind S.~465.}
med Conseqvens og Klarhed fremstillet
Forholdet imellem det Theoretiske og det Praktiske,
det Logiske og det Ethiske. „Kun Tænkningens Dannelse“,
siger han, „er Menneskets Dannelse, thi Tanken er paa
eengang det \emph{Menneskelige}, som adskiller os fra den
underordnede Natur, og det \emph{Guddommelige}, som for\-%
ener os med den os overordnede Verden. Maaskee vil
man indvende, at det Samme kan siges om \emph{Fölelsen};
det kan det ogsaa, men kun forsaavidt som denne er
gjennemtrængt af Tanken. Tro, Kjærlighed, Ære, Pligt
ere visselig Fölelser, som adskille os fra Naturen og
indlemme os i en höiere Verden; men hvad der giver
disse Fölelser saa stort et Værd, ere de i dem indeholdte
Tankebestemmelser, uden hvilke de ikke kunde existere
i Form af Fölelser. Hvis der f.\ Ex.\ ikke var en moralsk
Verden, hvori det Gode og det Onde vare adskilte ifölge
Tankens fornuftige Bestemmelser, saa kunde vi heller
ikke adskille dem i vor Fölelse; vi havde da ingen
% --- p. 68 ---
\opage{68}moralsk eller Pligtfölelse, ingen Kjærlighed til det Gode
og ingen Afsky for det Onde, thi alle disse Fölelser for\-%
udsætte en Magt, som bestemmer dem; denne Magt er
Tanken, som er den eneste absolute Bestemmelsesgrund.\,.\,.
% three close-set points, as printed
Staaer det altsaa fast, at Tanken baade er det Menneske\-%
lige og det Guddommelige i os, da er det ligesaa vist,
at Dannelsen af vor Tænkning er Dannelsen af os selv“.
% closing quote before the full stop, as printed
Naar Tænkningen paa denne Maade gjöres til Eet med
Mennesket, saa at Menneskets eiendommelige Væsen kun
bliver en Modification af den almene, objective Tænkning,
da fölger det af sig selv, at Villien trods sit subjective
Væsen ogsaa maa være en Modification af den objective
Tanke.

Men hvo seer nu ikke, at slige Synsmaader maae lede
os ind paa yderst farlige Afveie? Som Opium virker be\-%
dövende paa Nerverne, saaledes virker den objective
Tænkning, naar den skal anvendes efter ovenstaaende
Theorie, bedövende paa den praktiske Bevidsthed; og
ligesom Nydelsen af Opium i Begyndelsen kan erstatte
den tabte Virkelighed ved ubeskrivelige Henrykkelser og
paradisiske Drömmesyner, men, naar Drömmen er forbi,
efterlader frygtelige Virkninger, saaledes kan ogsaa en
Forgudelse af Theoriens objective Tænkning i Begyndel\-%
sen henrykke den Betragter, der seer Tilværelsens Idealer
i Speilet, men naar Betragtningen er tilende, staaer Virke\-%
ligheden kun endnu mere kantet, haard og trodsig mod Ide\-%
alitetens Ret; det er da forgjæves, at Tænkningen vil pröve
paa at bortforklare en Realitet, den ikke kan bære.

I Forhold til Virkeligheden og det virkelige Livs Op\-%
gaver maa den praktiske Erkjendelse absolut have större
Oprindelighed end den theoretiske. Vel maa Villien, for
at være Villie, oprindelig forudsætte Tanken; men deels
er det, vel at mærke, ikke den theoretiske, men den
praktiske Tænkning, Villien forudsætter, deels er Villien
med denne sin Forudsætten ligesaa oprindelig som Tanken.
Ligesom der i Tanken er Noget, der ikke kan afledes af
% --- p. 69 ---
\opage{69}Villien, efterdi Tanken som Tanke ikke er Villie, saa\-%
ledes er der ogsaa i Villien Noget, der ikke kan afledes
af Tanken, efterdi Villien som Villie ikke er Tanke.
Vistnok falde saadanne Kategorier som Valgfrihed, Valg,
Beslutning o.\ fl.\ ind under Tankens Belysning, forsaavidt
de klare sig i sædelig Erkjendelse; men i Principet ere
de ligefuldt Villieskategorier, deres Erkjendelse er væ\-%
sentlig praktisk. Det, der erkjendes theoretisk, kan ikke
erkjendes praktisk; det nytter ikke at beraabe sig paa,
at Indholdet, der erkjendes, er det samme, kun Erkjen\-%
delsesmaaden forskjellig: i enhver afgjörende Erkjendelse
bliver Indholdet bestemt ved Maaden, og Maaden ved
Indholdet. At Modet som Dyd kan erkjendes i Theorien,
maa indrömmes; men dersom det Mod, der erkjendes
theoretisk, virkelig i Eet og Alt var det samme, det selv\-%
samme som det, der erkjendes praktisk, hvad vilde saa
heraf fölge? Heraf vilde jo fölge, at en dygtig theoretisk
Tænker, uden at vove Noget, uden at udsætte sit Liv
for nogen virkelig Fare, maatte kunne tænke sig saaledes
ind i det sande Mod, at han i Virkeligheden med det
Samme tænkte det sande Mod ind i sig, og saaledes ved
objectiv Tænkning, hvor feig han for Resten var af Na\-%
tur, virkelig blev modig. Men at den, der i Virkelig\-%
heden er feig, om han end tænker nok saa theoretisk,
ligesaa lidt kan tilspeculere sig virkeligt Mod, som den,
der er uden Penge, om han end vilde tænke theoretisk
objectivt, til han revnede, kan tilspeculere sig virkelige
Penge, trænger ikke til Beviis.

Hvor dyb den Klöft er, som paa Grund af de ueens\-%
artede Principer danner Skilsmissen imellem den theo\-%
retiske og praktiske Erkjendelse, kan endvidere sees
deraf, at den absolute Sandhed selv maa adskille sig i
to Sandheder, en upersonlig, der eensidig svarer til
Tænkningen, og en personlig, der svarer til Villien. At
Tænkningens Idee særlig faaer Navn af Sandheden, me\-%
dens Villiesideen, ifölge en almindelig vedtagen Sprog\-%
% --- p. 70 ---
\opage{70}brug faaer Navn af det Gode, har i nyere Tid ikke
været tilstrækkeligt til at afværge det theoretiske Over\-%
greb, ved hvilket de to Ideer ere blevne sammensmeltede
saaledes, at det Gode er gaaet under i det Sande, lige\-%
som Villien i Tænkningen. Vel sætter Theorien en For\-%
skjel imellem Tanken og Villien, imellem det Sande og
det Gode, men da Forskjellen selv falder indenfor Theo\-%
rien, fölgelig er og bliver theoretisk, beholder Tænkningen
naturligviis bestandig Overmagten.

Hver Gang der spörges om Gyldigheden enten af
Ethikens eller Religionens Forudsætninger, mærker man
Theoriens Anmasselse. Valgfriheden f.\ Ex.\ kan ikke
have ubetinget Gyldighed; og hvorfor ikke? Fordi den
objective Tænkning, naar den skal begribe Friheden,
ikke kan Andet end nedsætte Vilkaarligheden til et de
overveiende Motiver underordnet Moment, saa at Fri\-%
heden forvandles til en i Reflexionen gjennemsigtig Nöd\-%
vendighed. Men Villiens Realitet fordrer jo dog en ube\-%
tinget Selvbestemmelse, det Godes Idee fordrer en op\-%
rindelig, i og ved sig selv gyldig Frihed. Ja, svarer
Theorien, lad Villien som faktisk Vilkaarlighed kun gjöre
sig saa lystig den vil, lad det Godes Idee i Praxis
spænde sine Fordringer saa höit den vil: naar Spörgs\-%
maalet er om Villiens og det Godes Realitet, saa gjæl\-%
der det om Sandheden; og naar det gjælder om Sand\-%
heden, saa er det dog Tænkningen, den almeengyldige,
i Viden sig udviklende Tænkning, der er og bliver överste
Dommer. Er det Aabenbaringens Gud, Himlens og Jor\-%
dens almægtige Skaber, hvorom Talen er, strax veed Theo\-%
rien, at Troesartiklens „almægtige Skaber“ umulig kan være
den sande Gud; og hvorfor ikke? Fordi Skabelsesacten
umulig lader sig tænke som en Vilkaarlighedsact, altsaa
ikke, hvad Artiklen dog aabenbart forudsætter, som en
ubetinget Frihedsact! Men Religionen, den aabenbarede
Religion, der forkynder sig som det absolute Gode, kan
nu een Gang ikke bestaae uden en saadan Forudsætning;
% --- p. 71 ---
\opage{71}Gud selv vilde jo ophöre at være den Hellige, den i
Sandhed Almægtige, dersom han i sit Forhold til Verden
ikke var den absolut frie Gud, den Gud, der skaber af
Intet. Religion hist, Religion her! hvad kan det hjælpe,
siger Theorien, at appellere enten til et aabenbaret Ords
umiddelbare Autoritet eller til Religionens populære Fore\-%
stillinger, naar Spörgsmaalet er om en evig Sandhed.
Kun for den objective Tanke er Sandheden objectiv; en
Gud kan i Folkephantasien være saa hellig han vil; er
han ikke Begrebets Gud, saa er han heller ikke nogen
sand Gud, er han ikke nogen sand Gud, saa har hans
Hellighed Intet at betyde.

Imod en saadan Vidensdespotisme maa der prote\-%
steres. At protestere imod Theoriernes Despotisme er
ingenlunde, hvad man fra et abstract theoretisk Stand\-%
punkt stundom har villet udgive det for, det Samme
som at krænke eller dog tilslöre den rige Sandhed, tvert\-%
imod, det er ikke alene i Religionens og Ethikens, men
i Sandhedens, ja i Theoriens eget Navn, paa höie Tid,
at der gjöres Indsigelser imod den objective Videns, den
theoretiske Erkjendelses utilbörlige Overgreb.

Den Protest, hvormed den praktiske Erkjendelse
afviser Theoriens Overgreb, er ikke et fortvivlet For\-%
sög af en Aandsretning, som i Mangel af objective
Grunde vil betjene sig af et Magtsprog, hidse Vilkaarlig\-%
heden og appellere til en blind Subjectivitet. Den prak\-%
tiske Erkjendelse er i dette Tilfælde saa langt fra at
voldtage den theoretiske, at den tvertimod begynder
med Friheden, den oprindelige Frihed, som Theorien,
naar den ret kommer til theoretisk Besindelse, selv maa
indrömme Subjectiviteten. Thi det i Subjectiviteten Aller\-%
subjectiveste er jo dog Selvet; men det Selviske i Selvet
er Selvbestemmen, er Villie. Seer nu den theoretiske
Erkjendelse sig nödt til at indrömme Villien og Friheden
samme Oprindelighed, som den indrömmer Tænkningen
og Reflexionen, da seer den sig ogsaa nödt til at ind\-%
% --- p. 72 ---
\opage{72}römme den subjective, personlige Sandhed, Villiens,
det Godes Idee, samme Oprindelighed, som den ind\-%
römmer Tænkningsideen, den upersonlige, objective Sand\-%
hed. „Og hvad“, spörger Theoretikeren, „kan den prak\-%
tiske Erkjendelse saa forlange mere?“ Den kan, svare vi,
forlange, at Theorien gjör Alvor af sine egne Indröm\-%
melser; den kan forlange, at den praktiske Idee i alt
Praktisk maa være ligesaa selvstændig, altsaa ligesaa
uafhængig af uvedkommende Theorier, som den theo\-%
retiske Idee i Alt, hvad der udelukkende angaaer
det Theoretiske, maa være fritagen for praktiske Inter\-%
essers forstyrrende Indblanding; den kan forlange Villies\-%
principet respecteret; den kan forlange, at Frihedens Mynt,
til Trods for objective, var det end dogmatisk-objective
Speculationer, bliver omsat fra Theoriens Navneværdi i
Charakterens Valuta, i Existensens Sölvværdi.

Paa disse, men ogsaa kun paa disse Vilkaar, vil
der kunne tilveiebringes sand Forstaaelse imellem Theorie
og Praxis. Er denne Forstaaelse först tilveiebragt, ikke
overfladisk, men fra Grunden af, da vil en Forsoning af
Tro og Viden, Religion og Videnskab, fölge som af sig
selv; thi Religionen er fra Först til Sidst praktisk, Viden\-%
skaben fra Först til Sidst theoretisk: Troen er et Villiens,
Viden et Tankens Ideal.

Det er en Selvfölge, at Tankens og Villiens Idealer
maae kunne indgaae i et ligesaa mangesidigt Vexelfor\-%
hold som Tanken og Villien selv; ethvert Forsög paa at
drage en endelig bestemt Grændselinie, der skulde ad\-%
skille den theoretiske Erkjendelse fra den praktiske, er
forsaavidt urimeligt og meningslöst. Men derfor er her
dog ligefuldt en Grændse, en principiel Forskjel, der
fremkommer derved, at alle Sjælekræfter i Theorien
underordnes Tænkningen, medens de i Praxis underordnes
Villien.

Vel kan der af enhver Livets Idee gives baade en
theoretisk og praktisk Erkjendelse; men hvor gjennem\-%
% --- p. 73 ---
\opage{73}gaaende Forskjellen er, er iöinefaldende ved et Blik paa
Gudsideen. Det er nödvendigt at sige det resolut: Tan\-%
kens Gud, hvis Erkjendelse er theoretisk, kan ikke være
Troens Gud, thi af Troens Gud gives der kun en prak\-%
tisk Erkjendelse. Stærkere kan Modsætningen mellem
Tro og Viden, Religion og Videnskab, praktisk og theo\-%
retisk Erkjendelse ikke udtales, end naar det udtrykke\-%
lig siges, at de have hver sin Gud.

„Og dette“, spörger man, „skal være den gode For\-%
staaelse“? Vi svare: Ja, det er den bedste Forstaaelse,
der lader sig tænke. Naar Videnskaben tilfulde indseer,
at den ligesaa lidt kan erkjende Troens Gud, som Reli\-%
gionen kan erkjende Videns Gud, idet Religionen fra sin
praktiske Erkjendelses Standpunkt indseer det Samme,
hvad saa? „Ja, saa“, udbryder man, „maae de först for
Alvor komme i Strid med hinanden, thi enhver af dem
vil naturligviis antage sin Gud, ikke blot for en sand
Gud, men for den eneste sande Gud“! Slige Expectora\-%
tioner, som, hvad Theorien angaaer, mere charakterisere
de livlige end de grundige Hoveder, grunde sig aaben\-%
bart paa et overfladisk Raisonnement.

Naar det siges, at Tro og Viden umulig kunne have
den samme Gud, saa er Meningen naturligviis ikke den,
at Troens Gud skulde være en Anden end Videns i
samme Forstand, som Keiseren af China er en Anden
end Keiseren af Brasilien; men Meningen er, at det Gud\-%
dommelige, der tænkes, ikke er Gjenstand for Tro, og
at Gud, forsaavidt han er Gjenstand for Tro, ikke er
og ifölge Sagens Natur ikke kan være Object for be\-%
gribende Tænkning.

At Aandsvidenskaben, forsaavidt den forstaaer sig
selv, maa være villig til at gaae ind paa den foreslaaede
Adskillelse, er naturligt; thi aldrig har nogen Philosophie
endnu udviklet Gudsbegrebet saaledes, at den i Begrebet
erkjendte Gud kunde blive Gjenstand for Tro. Den, der
paastaaer, at Nogen virkelig kan troe enten paa Aristo\-%
% --- p. 74 ---
\opage{74}teles's Gud eller paa Leibnitz's Gud eller paa Hegels
Gud, er tankelös og veed ikke, hvad han siger. Men
at Religionen fra sin Side heller ikke kan have Noget imod
Adskillelsen, ligger ligeledes i Sagens Natur, og er paa
utvetydig Maade indrömmet af Apostelen Paulus.

% "19-22": a short hyphen-like dash, as printed; contrast the long dash in
% "(22—32)" below.
Idet Paulus (Rom.\ 1, 19-22) taler om det Eien\-%
dommelige i Hedningenes Gudserkjendelse, lægger han
aabenbart Vægt paa det i Guds Væren, der lader sig
tænke, paa „det, som man kan vide om Gud“ (τὸ γνωστὸν
τοῦ θεοῦ), paa Guds „usynlige Væsen, hans evige Kraft
og Guddommelighed“ (ἣ τε ἀΐδιος αυτοῦ δύναμις καὶ
θειότης), „der beskues og forstaaes“ (νοούμενα καθορᾶται)
% GREEK (Rom. 1:19--20), set in a sloping Greek fount; typed here with the
% accents and breathings the page shows. The print's variant forms (the
% looped "6"-shaped sigma in γνωστὸν, the curly theta in θεοῦ, θειότης,
% καθορᾶται, the open kappa in καὶ, καθορᾶται) are normalised to σ θ κ.
% AS PRINTED: "ἣ" with rough breathing and GRAVE (the NT reads ἥ, acute);
% "αυτοῦ" with NO breathing (NT αὐτοῦ) — both checked at 1200 ppi, printer's
% readings, not corrected.
af hans Gjerninger. Det er altsaa i strengere Forstand
ikke Villiens, men Tankens Gud, der har givet sig til\-%
kjende for Hedningene. Det er den i Naturriget aaben\-%
barede, Væsenhedens, Grundens, Magtens, Lovenes og
den rationelle Nödvendigheds-Guddom, hvis intellectuelle
Lys har kunnet skinne for Naturfolkene. Men at en Guds\-%
erkjendelse, som paa disse Vilkaar maa blive mere theo\-%
retisk end praktisk, altsaa mere sigte paa Væsen end
paa Villie, kun kan afgive et svagt Grundlag for det
Ethiske, at Gudstheorien ikke har kunnet sikkre Prin\-%
ciperne for det sædelige Liv, anskueliggjöre de fölgende
Vers (22---32) paa en saa indtrængende Maade, at det
af Fordærvelsen udkastede Billede maa forfærde os.

Hvor forskjellig den hedenske Gudserkjendelse er
fra den jödiske og den christelige, kan sees af Romer\-%
brevets 9de Capitel. Thi hvad er det dog for en Gud,
hvorom Paulus taler i dette Capitel? Det er ikke Tan\-%
kens Gud, hvis Frihed dæmrer i overgribende Nödven\-%
dighed; nei, det er Frihedens Gud, hvis Væsen er ube\-%
tinget Villie. Vel har den for Hedningene „aabenbarede“
Gud ogsaa en Villie, men det er kun en Tankevillie, en
Villie, der er Moment for det med den objective Sand\-%
hed identiske Gode (Deus vult, quia bonum est). Israels
Gud er derimod lutter Villie, hans Tanker Villiestanker,
% --- p. 75 ---
\opage{75}hans Slutninger Raadslutninger, hans Udvælgelser idel
Vilkaarlighed. Thi at Aabenbaringens Gud fremstilles
som Vilkaarlighedens Gud, vil jo ingenlunde sige, at
Guds ubetingede Villie skulde drive Spil enten med Nöd\-%
vendigheden eller med Sandheden eller med Godheden;
tvertimod! Al Grund, al Sandhed og Godhed har netop
% "11": the first 1 is an ink blot at 600 ppi; the second is clean. Read 11
% (Rom. 11:32--33 fits the sense); the dash is shorter than an em rule.
(Rom.\ 11, 32--33) sit Midtpunkt, sit Udgangspunkt og
Hvilepunkt i den guddommelige Villie (bonum est, quia
Deus vult).

Hvor umuligt det er at udvikle nogen objectiv Theorie
om Villiens Gud, kan da Enhver, som vil eftersee Romer\-%
brevet, ved nogen Eftertanke sige sig selv; og skulde
han föle Trang til yderligere Beviis, da ville Skriftfor\-%
tolkningens theologiske Kvaklerier, Dogmatikens specula\-%
tive Selvmodsigelser og Prædestinationstheoriernes Abra\-%
cadabra levere Beviser i Mængde. Men det Samme, der
gjör, at Vilkaarlighedens Gud ikke kan være Gjenstand for
theoretisk Erkjendelse, det Selvsamme gjör ogsaa, at
Villiens aabenbarede Gud er og maa være Gjenstand for
Tro. Som Gjenstand for Tro er Villiens Gud oprindelig
Underets Gud; thi hvad er vel Underet andet end den
guddommelige Villies ubetingede Magt over alle Betin\-%
gelser. Den, der sætter sit Liv i en almindelig Tænken,
faaer maaskee en Theorie om Livet, men for Livet selv
bliver han bedragen; den, der sætter Tanken ind paa
praktisk Erkjendelse af en hellig Villie, gaaer maaskee
glip af Theorien, men han vinder Livet, Aandens evige
Liv, thi Aandens Liv er i Villien.

Hvor inderlig Troen og den praktiske Erkjendelse
hænge sammen med Villien, kan sees paa Lydigheden.
At troe paa Gud er at lyde Gud, er at gjöre Guds Villie;
men at lyde Gud er igjen det Samme som at troe paa
Gud, thi Villiens förste Bud er Troen. Hvor inderlig
Troen og Lydigheden svare til hinanden, kan sees paa
de to meest fremragende Phænomener af religiös Erkjen\-%
delse, som vor Tid har at opvise, Grundtvig og S.\ Kierke\-%
% --- p. 76 ---
\opage{76}gaard: \emph{Troen er væsentlig Lydighed}, dette er
i Korthed Kierkegaards Forklaring; \emph{Lydighed} er
\emph{væsentlig Tro, Tro paa det Livets Ord, der er
Menneskenes Lys}, dette er i en Hovedsum Grundt\-%
vigs Forklaring. De to Forklaringer ere saa langt fra
at staae i indbyrdes Strid, at de meget mere fuldstændig\-%
gjöre, belyse og stadfæste hinanden.

Skulde det endnu virkelig, hvad Nogle paastaae,
være den ene og samme Bevidsthed umuligt at aner\-%
kjende baade det Praktiske og det Theoretiske, baade
Troen og Viden, da vilde det i al Fald blive det Theo\-%
retiske, ikke det Praktiske, der maatte opgives. Exi\-%
stensen er praktisk; en Bevidsthed kan existere uden at
være theoretisk; men en Bevidsthed kan ikke være theo\-%
retisk uden at existere. Det er altsaa muligt at sætte
Existensen överst og beholde Theorien; men at aflede
Existensen af Theorien, at sætte Theorien överst og be\-%
holde Existensen, er umuligt. Existensidealet er Villie,
thi Existensen er praktisk; Tankeidealet er Viden, thi
Tanken, den objective Tanke, er theoretisk: at sætte
Existensidealet överst og beholde Tankeidealet er mu\-%
ligt, men at sætte Tankeidealet överst og fyldestgjöre
Existensidealet er umuligt; med andre Ord: at sætte
Troen överst og beholde Videnskaben er muligt; men
at sætte Videnskaben överst og beholde Troen er umuligt.

% "Nogen": the o carries a single grave-like stroke at 600 ppi, not a
% diaeresis -- a damaged or foreign sort; read as Nogen.
Skulde Nogen imidlertid være saa livlig i Hovedet
og saa let i Tankerne, at det virkelig löber rundt for ham,
saa at han ikke kan finde Rede i, hvad der skal staae
överst og hvad der skal staae nederst, Troen eller Viden,
men af lutter Livlighed forvexler Överst med Nederst
og Nederst med Överst, lad ham forsöge at simplificere
Forholdene ved at kaste Troen, Villiesidealet, Charak\-%
teerprincipet, bort, og leve paa den rene Theorie: han
undgaaer dog ikke Fordoblingen. Fordoblingen af Theo\-%
retisk og Praktisk lægges da ind i Theorien. For at be\-%
vare Selvbevidsthedens Eenhed, skulde et Menneske alt\-%
% --- p. 77 ---
\opage{77}saa have sin Theorie i det Praktiske, sin Praxis i det
Theoretiske, sin Tænken i Villen og sin Villen i Tænken;
% "Villen" (twice), not "Villien": the verbal noun, paired with "Tænken";
% as printed.
men ved at löse en saadan Opgave skal vist ingen Livlig
rose sig af at have löst Tilværelsens Gaade. Ved at
fordoble sig praktisk indenfor det Theoretiske bliver man
med al mulig Livlighed i livlig Fordobling --- ikke en
ethisk Charakteer, men et intellectuelt Amphibium.

% END OF ARTICLE. Nothing closes it but a short centred rule below the last
% line (no signature, no date); the rest of p. 77 is blank. The Nekrologer
% begin on the next page (PDF 652, printed 78, folio suppressed).
\begin{center}\rule{0.2\linewidth}{0.4pt}\end{center}

\end{document}
